% META: {"id": "TCOMPL-INEG-EV-A", "chapitre": "TCOMPL-INEGALITES", "type_objet": "evaluation", "version": "A", "duree_min": 45, "bareme_total": 20, "capacites_codes": ["C2", "C4", "C5"], "competences": ["calculer", "raisonner"], "mode_creation": "ex_nihilo", "sources_inspiration": [], "status": "generated"}
% BEGIN-VERIFY
% from sympy import symbols, diff, integrate, Rational
% x = symbols('x')
% L = (3*x**2 + x) / 4
% assert L.subs(x, 0) == 0
% assert L.subs(x, 1) == 1
% Lp = diff(L, x)
% assert Lp.subs(x, 0) > 0
% assert Lp.subs(x, 1) > 0
% Lpp = diff(L, x, 2)
% assert Lpp == Rational(3, 2)
% I = integrate(L, (x, 0, 1))
% G = 1 - 2*I
% assert G == Rational(1, 4)
% END-VERIFY

%---------------------------------------------------------------
% Duree : 45 minutes — Bareme : 20 points
%---------------------------------------------------------------

\begin{evaluation}{TCOMPL-INEG-EV-A}{45}

\begin{exercice}{TCOMPL-INEG-EV-A-EX1}{2}{40}
\textbf{Exercice unique} (20 points) --- \textit{Capacités C2, C4, C5}

\medskip

On modélise une courbe de Lorenz par $L(x) = \dfrac{3x^2+x}{4}$ sur $[0,1]$.
\begin{enumerate}
  \item (4 pts) Vérifier que $L(0)=0$ et $L(1)=1$.
  \item (4 pts) Calculer $L'(x)$ et vérifier que $L$ est croissante sur $[0,1]$ (indication : étudier le signe de $L'$).
  \item (4 pts) Calculer $L''(x)$ et en déduire que $L$ est convexe.
  \item (8 pts) Calculer l'indice de Gini associé à ce modèle.
\end{enumerate}
\end{exercice}

\end{evaluation}
